\documentclass[landscape, 8pt]{extarticle}

\usepackage{../../Key-Files/preamble}
\usepackage[separate]{../../Key-Files/thmboxes_v5}

\DeclareMathOperator{\epi}{epi}


\begin{document}
	
%\fontsize{4pt}{5pt} \selectfont %Smallest legible font size when printing with UoE Main Library printers, however some resolution issues 	for extremely small charactes such as subscripts inside fractions. Only use if you *really* need the extra space
\fontsize{5pt}{6pt} \selectfont
	
%Reduce horizontal padding in 'tabular' environment, make rows taller for legibility
\setlength\tabcolsep{4.5pt}
\setlength\extrarowheight{2pt}
	
\setlength{\abovedisplayskip}{0pt}
\setlength{\belowdisplayskip}{0pt}
\setlength{\abovedisplayshortskip}{0pt}
\setlength{\belowdisplayshortskip}{0pt}

%%%%%%%%%%%%%%%%%%%%%%%%%%%%%%%%%%%%%%%%%%%%%%%%%%%%%%%%%%%%%%%%%%%%%%%%%%%%%%%%%%%%%%%%%%%%
%%%%%%%%%%%%%%%%%%%%%%%%%%%%%%%%%% END PREAMBLE %%%%%%%%%%%%%%%%%%%%%%%%%%%%%%%%%%%%%%%%%%%%
%%%%%%%%%%%%%%%%%%%%%%%%%%%%%%%%%%%%%%%%%%%%%%%%%%%%%%%%%%%%%%%%%%%%%%%%%%%%%%%%%%%%%%%%%%%%	
\begin{multicols*}{4}
	\raggedcolumns 
	%Original template by Leon, modified by Louis
	\vspace{-5pt}
	\subsection*{FuO Exam Sheet 2026/7 - }
	\vspace{-5pt}
	
%%%%%%%%%%%%%%%%%%%%%%%%%%%%%%%%%%%%%%%%%%%%%%%%%%%%%%%%%%%%%%%%%%%%%%%%%%%%%%%%%%%%%%%%%%%%
%%%%%%%%%%%%%%%%%%%%%%%%%%%%%%%%%%%%%%%% BASICS %%%%%%%%%%%%%%%%%%%%%%%%%%%%%%%%%%%%%%%%%%%%
%%%%%%%%%%%%%%%%%%%%%%%%%%%%%%%%%%%%%%%%%%%%%%%%%%%%%%%%%%%%%%%%%%%%%%%%%%%%%%%%%%%%%%%%%%%%

{\footnotesize \textbf{\underline{Basics}}}
		
		
		\begin{dfn}[Ingredients]{dfn:ingredients}{ }
			An optimisation problem has four ingredients:
			\begin{itemize}
				\item \textbf{Decision variables} $\mathbf{x}\in \mathbb{R}^n$: Set of controllable inputs we optimise over
				\item \textbf{Parameters}: Set of fixed uncontrollable inputs 
				\item \textbf{Feasible Region} $\mathcal{S} \subseteq\mathbb{R}^n$ Description of the set of candidate solns
				\item \textbf{Objective function} $f: \mathbb{R}^n \to \mathbb{R}$: Measure of how 'good' a candidate soln is
			\end{itemize}
			And thus we express an optimisation problem as $(P) \quad \min \{f(\mathbf{x}) : \mathbf{x}\in \mathcal{S}\}$ or equivalently $(P) \quad \min f(\mathbf{x})$ subject to $\mathbf{x}\in \mathcal{S}$.
			\begin{itemize}
				\item Any max problem can be made a min problem by negating $f$
				\item Recall a func $g: \mathbb{R}^n \to \mathbb{R}$ is \textbf{linear} if $g(\alpha \mathbf{x} + \beta \mathbf{y}) = \alpha g(\mathbf{x}) + \beta g(\mathbf{y}) \; \forall \; \mathbf{x, y}\in \mathbb{R}^n$
				\begin{itemize}
					\item $g: \mathbb{R}^n \to \mathbb{R}$ is linear iff $\exists \mathbf{a}\in \mathbb{R}^n$ s.t. $g(\mathbf{x}) = \mathbf{a}^T \mathbf{x} \; \forall \; \mathbf{x}\in \mathbb{R}^n$
					\item Every linear function $g$ has $g(\mathbf{0}) = 0$
				\end{itemize}
			\end{itemize}
			
		\end{dfn}
		
		
		\begin{dfn}[Terminology]{dfn:terminology}{ }
			\begin{itemize}
				\item $\mathcal{S}$ is the \textbf{feasible region} or \textbf{feasible set}
				\item Any soln $\mathbf{\hat{x}}\in \mathcal{S}$ is a \textbf{feasible solution}
				\item $\mathbf{x}^*\in \mathcal{S}$ is an \textbf{optimal solution} if $f(\mathbf{x}^*) \le f(\mathbf{\hat{x}}) \forall \mathbf{\hat{x}}\in \mathcal{S}$
				\item $\mathcal{S}^*$ is the \textbf{set of all optimal solutions}
				\item $z^*\in \mathbb{R} \cup \{-\infty\} \cup \{\infty\}$ is the \textbf{optimal value} if $z^* \le f(\mathbf{\hat{x}}) \forall \mathbf{\hat{x}}\in \mathcal{S}$ and $\exists$ a sequence $\mathbf{x}^{(k)}\in \mathcal{S}$, $k\in \mathbb{N}$, s.t. $f(\mathbf{x}^{(k)}) \to z^*$ as $k \to \infty$
				\item If $\mathcal{S}^* \ne \emptyset$ then $z^* = f(\mathbf{x}^*) \forall \mathbf{x}^*\in \mathcal{S}^*$ and $z^*$ is \textbf{attained}
				\item $(P)$ is \textbf{unbounded} if $\exists$ sequence $\mathbf{x}^{(k)}\in \mathcal{S}, k\in \mathbb{N}$ s.t. $f(\mathbf{x}^{(k)}) \to -\infty$ as $k \to \infty$. Here $\mathcal{S}^* = \emptyset$ and $z^* := -\infty$
				\item If $\mathcal{S} = \emptyset$ then $(P)$ is \textbf{infeasible} and $z^* := +\infty$
				\item If $(P)$ is not infeasible or unbdd, $-\infty < z^* < \infty$, but $\mathcal{S}^*$ may still $= \emptyset$
			\end{itemize}
			
			
		\end{dfn}
		
		
		\begin{mdfn}[Constrained Optimisation]{dfn:constrainedopt}{ }
			If $\mathcal{S} = \mathbb{R}^n$, $(P)$ is \textbf{unconstrained}, else $\mathcal{S} \subset \mathbb{R}^n$ and $(P)$ is \textbf{constrained}
			\begin{itemize}
				\item We typically express $\mathcal{S}$ with functional relations:
					\begin{itemize}
						\item $\mathcal{S} = \{\mathbf{x}\in \mathbb{R}^n : g_i(\mathbf{x}) \ge b_i, i\in M_1;\; \ell_i(\mathbf{x}) \le b_i, i\in M_2; \; h_i(\mathbf{x})=b_i, i\in M_3\}$
						\item $M_1, M_2, M_3$ are finite (maybe empty) index sets
						\item $g_i/\ell_i/h_i: \mathbb{R}^n \to \mathbb{R}$ for $i\in M_1/M_2/M_3$ and $b_i\in \mathbb{R} \; \forall \; i\in M_1 \cup M_2 \cup M_3$
					\end{itemize}
				\item We prefer $\ge/\le$ to $</>$ as $\mathcal{S}$ may not be closed with $</>$ inequalities\\
				\item $(P)$ is a \textbf{linear programming problem} if all of $f, g_i, \ell_i, h_i$ are linear
				\sbsadaptr{
					LPPs can be rephrased as: \\
					with $\mathbf{c} \in \mathbb{R}^n$ and $\mathbf{a}^{(i)} \in \mathbb{R}^n, b_i\in \mathbb{R}$ for all $i\in M_1 \cup M_2 \cup M_3$
				}{
				$ \begin{array}{ccc}
					(P) \;\;& \min & \mathbf{c}^T\mathbf{x}\\
					&\text{s.t.} & (\mathbf{a}^{(i)})^T\mathbf{x} \ge b_i, \; i \in M_1\\
					&&(\mathbf{a}^{(i)})^T\mathbf{x} \le b_i, \; i \in M_2\\
					&&(\mathbf{a}^{(i)})^T\mathbf{x} = b_i, \; i \in M_3
				\end{array} $
				}
				\item The feasible region of any LPP is a polyhedron
			\end{itemize}
			
		\end{mdfn}
		
		
		\begin{dfn}[Convex Sets]{dfn:convexsets}{ }
			A set $\mathcal{C} \subseteq \mathbb{R}^n$ is \textbf{convex} if for all $\mathbf{x, y}\in \mathcal{C}$ and $\lambda\in [0, 1]$, we have $\lambda \mathbf{x} + (1-\lambda)\mathbf{y}\in \mathcal{C}$, otherwise it is \textbf{nonconvex}
			\begin{itemize}
				\item $\emptyset, \mathbb{R}^n$, and single-element sets are convex
				\item Convexity is preserved under intersection but not necessarily union
				\item LPPs are convex optimisation problems, so $\mathcal{S}$ and $\mathcal{S}^*$ are convex
			\end{itemize}
			
		\end{dfn}
		
		
		\begin{dfn}[Convex/Concave Functions]{dfn:convexfunctions}{ }
			$f: \mathbb{R}^N \to \mathbb{R}$ is \textbf{convex} if $\forall \; \mathbf{x, y}\in \mathbb{R}^n$ and $\lambda\in [0, 1]$, $$f(\lambda \mathbf{x} + (1-\lambda)\mathbf{y}) \le \lambda f(\mathbf{x}) + (1-\lambda)f(\mathbf{y}).$$
			\begin{itemize}
				\item Any linear function is convex and concave
				\item Intuition - convex if line between any two points on the graph of $f$ do not intersect the graph
				\item The \textbf{epigraph} of $f$, $\epi(f)$ is the region above the graph of $f$
				\begin{itemize}
					\item $f$ is a convex func iff $\epi(f)$ is a convex set
				\end{itemize}
			\end{itemize}
			$f$ is \textbf{concave} if the inequality is reversed. $f$ is concave iff $-f$ is convex.
		\end{dfn}
		
		
		\begin{dfn}[Level Sets]{dfn:levelsets}{ }
			For $f: \mathbb{R}^n \to \mathbb{R}$ and $\alpha\in \mathbb{R}$, the \textbf{level set of $f$} is $\mathcal{L}_{\alpha}(f) = \{x\in \mathbb{R}^n : f(x) = \alpha\}$
			\begin{itemize}
				\item $\mathcal{L}_\alpha(f)$ is the inverse image of $\alpha$ under $f$
				\item The union across all $a\in\mathbb{R}$ is equal to $\mathbb{R}^n$
				\item $\mathcal{L}_\alpha(f) = \emptyset$ iff $\alpha$ not in the range of $f$
				\item The sublevel (superlevel) set of $f$ is $\mathcal{L}_\alpha^-$ $(\mathcal{L}_\alpha^+)= \{x\in \mathbb{R}^n : f(x) \le (\ge) \alpha\}$
				\begin{itemize}
					\item Sublevel sets of convex functions are convex for all $\alpha\in \mathbb{R}$
					\item Superlevel sets of concave functions are convex for all $\alpha\in \mathbb{R}$
				\end{itemize}
			\end{itemize}
		\end{dfn}
		
		
		\begin{dfn}[Hyperplanes]{dfn:hyperplanes}{ }
			A \textbf{hyperplane} $\mathcal{H}$ is a level set of a linear function $f(\mathbf{x})$ for $\alpha\in \mathbb{R}$, can be written as $\mathcal{H} = \{\mathbf{x}\in \mathbb{R}^n : \mathbf{a}^T\mathbf{x} = \alpha\}$ where $f(\mathbf{x}) = \mathbf{a}^T \mathbf{x}$ for some $\mathbf{a}\in \mathbb{R}^n$.
			\begin{itemize}
				\item For linear $f:\mathbb{R}^n \to \mathbb{R}$, $\mathcal{L}_\alpha(f)$ is a hyplane with $n-1$ \textbf{degrees of freedom}
				\item Every hyperplane is a convex set
			\end{itemize}
			
		\end{dfn}
		
		
		\begin{dfn}[Halfspaces]{dfn:halfspaces}{ }
			A \textbf{halfspace} is a sublevel ($\mathcal{H}^-$) or superlevel ($\mathcal{H}^+$) set of a linear $f$
			\begin{itemize}
				\item Every halfspace is convex
				\item $\mathcal{H}$ is the boundary of $\mathcal{H}^-$ and $\mathcal{H}^+$, and $\mathcal{H}^- \cap \mathcal{H}^+ = \mathcal{H}$
				\item $\mathbf{a}\in \mathbb{R}$ is perpendicular to $\mathcal{H}$ ($\forall \mathbf{x}_1, \mathbf{x}_2\in \mathcal{H}$, $\mathbf{c}^T(\mathbf{x_1} - \mathbf{x}_2) = \alpha - \alpha = 0)$
			\end{itemize}
			
		\end{dfn}
		
		
		\begin{dfn}[Polyhedra \& Polytopes]{dfn:polyhedranpolytopes}{ }
			A set $\mathcal{P} \subseteq\mathbb{R}^n$ is a \textbf{polyhedron} if it is given by the intersection of a finite number of hyperplanes and halfspaces.
			\begin{itemize}
				\item All polyhedra are convex (they are the intersection of convex sets)
				\item $\mathcal{S} \subseteq \mathbb{R}^n$ is \textbf{bounded} if $\exists K\in \mathbb{R}$ s.t. $K \ge \left|x_j\right|$ for all components $x_j$ of $\mathbf{x}\in \mathcal{S}$
				\item A \textbf{polytope} is a bounded polyhedron
			\end{itemize}
			
		\end{dfn}
		
		
		


%%%%%%%%%%%%%%%%%%%%%%%%%%%%%%%%%%%%%%%%%%%%%%%%%%%%%%%%%%%%%%%%%%%%%%%%%%%%%%%%%%%%%%%%%%%%
%%%%%%%%%%%%%%%%%%%%%%%%%%%%%%%%%%%%%%%%%%%%% VERTICES %%%%%%%%%%%%%%%%%%%%%%%%%%%%%%%%%%%%%
%%%%%%%%%%%%%%%%%%%%%%%%%%%%%%%%%%%%%%%%%%%%%%%%%%%%%%%%%%%%%%%%%%%%%%%%%%%%%%%%%%%%%%%%%%%%

{\footnotesize \textbf{\underline{Vertices}}}
		
		
		
		\begin{thm}[Relating Halfspaces To Constrained Optimisation]{thm:relatehalfspacestoproblem}{ }
			In the generic constrained optimisation problem (P), the feasible region $\mathcal{S} \subseteq \mathbb{R}^n$ is 
			$$
			\mathcal{S} = \{x\in \mathbb{R}^n : g_i(x) \ge b_i, i\in M_1;\; \ell_i(x) \le b_i, i\in M_2; \; h_i(x)=b_i, i\in M_3\}
			$$$$
			\mathcal{S} = \bigg(\bigcap_{i\in M_1} \mathcal{L}^+_{b_i}(g_i)\bigg) \cap \bigg(\bigcap_{i\in M_2} \mathcal{L}_{b_i}^-(\ell_i)\bigg) \cap \bigg(\bigcap_{i\in M_3}\mathcal{L}_{b_i}(h_i)\bigg)
			$$
			
		\end{thm}
		
		
		\begin{mdfn}[Convex Optimisation Problems]{dfn:convexoptimisation}{ }
			The generic (P) is a \textbf{convex optimisation problem} if: $f$ is convex; all $g_i$ are concave; all $\ell_i$ are convex; and all $h_i$ are linear.\\
			Otherwise it is a \textbf{nonconvex optimisation problem}
			\begin{itemize}
				\item Set of LPPs $\subset$ set of convex optimisation problems
				\item For any LPP $\mathcal{S}$ and $\mathcal{S}^*$ are convex sets
			\end{itemize}
			
		\end{mdfn}
		
		
	{\tiny\textbf{\underline{Vertices Of Convex Sets}}}
		
		\begin{mdfn}[A Vertex Of A Convex Set]{dfn:vertexofconvexset}{ }
			For $\mathcal{C} \subseteq\mathbb{R}^n$ a convex set, $\mathbf{\hat{x}}\in \mathcal{C}$ is a \textbf{vertex of $\mathcal{C}$} if $\exists$ a hyperplane $\mathcal{H}$ with halfspace $\mathcal{H}^+$ such that both $\mathcal{C} \cap \mathcal{H} = \{\mathbf{\hat{x}}\}$; and $\mathcal{C} \subseteq \mathcal{H}^+$
			Such a hyperplane is a \textbf{supporting hyperplane} of $\mathcal{C}$\\
			Intuition - supporting hyperplanes are tangents to the convex set. They can intersect multiple vertices, if they do then they intersect the edge between those vertices.
			
		\end{mdfn}
		
		
		\begin{mthm}[Vertices As Optimal Solutions]{thm:verticesasoptima}{ }
			$\mathbf{\hat{x}}\in$ convex $\mathcal{C}$ is a vertex of $\mathcal{C}$ iff there exists linear $\ell(\mathbf{x}) = \mathbf{c}^T\mathbf{x}$ for $\mathbf{c}\in \mathbb{R}^n\backslash\mathbf{0}$ such that $\mathbf{\hat{x}}$ is the \textbf{unique optimal solution} to $\min \; \mathbf{c}^T \mathbf{x}$ s.t. $\mathbf{x}\in \mathcal{C}$.
			
		\end{mthm}
		
		
		{\tiny\textbf{\underline{Characterising Vertices}}}
		
		
		\begin{mdfn}[Active (aka Binding) Constraints]{dfn:activeconstraints}{ }
			For $\mathcal{P}$ a polyhedron and $\mathbf{\hat{x}}\in \mathbb{R}^n$, the set of \textbf{active}  constraints at $\mathbf{\hat{x}}$ is $$
			I(\mathbf{\hat{x}}) = \{i\in M_1 \cup M_2 \cup M_3 : (\mathbf{a}^i)^T\mathbf{\hat{x}} = b_i\}
			$$
			Note $M_3 \subseteq I(\mathbf{\hat{x}})$ for all $\mathbf{\hat{x}}\in \mathcal{P}$
			
		\end{mdfn}
		
		
		\begin{mdfn}[Basic \& Basic Feasible Solutions]{dfn:basicnbasicfeasiblesolns}{ }
			For $\mathcal{P} \subseteq \mathbb{R}^n$ a polyhedron and $\mathbf{\hat{x}}\in \mathbb{R}^n$, $\mathbf{\hat{x}}$ is a \textbf{basic solution} if both:\\
			\begin{enumerate}[label=(\roman*)]
				\item All equality constraints are active at $\mathbf{\hat{x}}$ (i.e. $M_3 \subseteq I(\mathbf{\hat{x}})$)
				\item $\{\mathbf{a}^i : i \in I(\mathbf{\hat{x}})\} \subset \mathbb{R}^n$ contains $n$ lin-ind vectors (ie spans $\mathbb{R}^n$)
			\end{enumerate} 
			$\mathbf{\hat{x}}$ is a \textbf{basic feasible solution} if it is basic and feasible
			
			\begin{itemize}
				\item Consider matrix $A(\mathbf{\hat{x}})$ with rows $(\mathbf{a}^{(i)})^T$ for $i\in I(\mathbf{\hat{x}})$. The no. of rows of $A$ is the number of active constraints at $\mathbf{\hat{x}}$, $\left|I(\mathbf{\hat{x}})\right|$. Thus $A$ has rank $n$ so the only vector in its null space is $\mathbf{0}\in \mathbb{R}^n$
				\item Equivalently $(\mathbf{a}^{(i)})^T \mathbf{d} = \mathbf{0}$ for $i\in I(\mathbf{\hat{x}})$ iff $\mathbf{d} = \mathbf{0}\in \mathbb{R}^n$
			\end{itemize}
			
		\end{mdfn}
		
		
		\begin{thm}[Basic Feasible iff A Vertex]{thm:basicfeasibleiffvertex}{ }
			For polyhedron $\mathcal{P} \subseteq\mathbb{R}^n$, $\mathbf{\hat{x}}\in \mathcal{P}$ is a bas-feas soln of $\mathcal{P}$ iff it is a vertex of $\mathcal{P}$
			\vspace{-7pt}
			\begin{proof}
				$\implies$: Need $\mathbf{a}\in \mathbb{R}^n\backslash \mathbf{0}$ and $\alpha\in \mathbb{R}$ s.t. $\mathbf{a}^T\mathbf{\hat{x}} = \alpha$ and $\mathbf{a}^T\mathbf{x} > \alpha$ for $\mathbf{x}\in \mathcal{P}\backslash\{\mathbf{\hat{x}}\}$. Take $I \subseteq I(\mathbf{\hat{x}})$ s.t. $\{\mathbf{a}^{(i)} : i\in I\}$ is a basis for $\mathbb{R}^n$. \\
				If $\mathbf{a} = \sum_{i\in I \cap (M_1 \cup M_3)} \mathbf{a}^{(i)}  + \sum_{i\in I \cap M_2}-\mathbf{a}^{(i)}$ then $\mathbf{a} \ne \mathbf{0}$ and \\
				$\alpha = \mathbf{a}^T\mathbf{\hat{x}} =  \sum_{i\in I \cap (M_1 \cup M_3)} (\mathbf{a}^{(i)})^T\mathbf{\hat{x}}  + \sum_{i\in I \cap M_2}-(\mathbf{a}^{(i)})^T\mathbf{\hat{x}} =  \sum_{i\in I \cap (M_1 \cup M_3)} b_i  + \sum_{i\in I \cap M_2}-b_i$. If $\mathcal{H} = \{x\in \mathbb{R}^n: \mathbf{a}^T\mathbf{x} = \alpha\}$ and $\mathcal{H}^+ = \{x\in \mathbb{R}^n : \mathbf{a}^T \mathbf{x} \ge \alpha\}$ then $\mathbf{\hat{x}}\in \mathcal{P} \cap \mathcal{H}$. For $\mathbf{x}\in \mathcal{P}$, $\mathbf{a}^T \mathbf{x} \ge \mathbf{a}^T\mathbf{\hat{x}} = \alpha$ by individual bounds of $b_i$ so $\mathcal{P} \subseteq \mathcal{H}^+$; and for any $\mathbf{x}\in \mathcal{P} \cap \mathcal{H}$, $(\mathbf{a}^{(i)})^T\mathbf{x} = b_i$ for $i\in I$, so $(\mathbf{a}^{(i)})^T(\mathbf{x} - \mathbf{\hat{x}}) = \mathbf{0}$ for $ i\in I$. Since $\{\mathbf{a}^{(i)} : i\in I\}$ is a basis, $\mathbf{x} - \mathbf{\hat{x}} = \mathbf{0} \implies \mathcal{P} \cap \mathcal{H} = \{\mathbf{\hat{x}}\}$. Thus $\mathbf{\hat{x}}$ is a vertex of $\mathcal{P}$.\\
				$\impliedby$: Know $\exists$ $\mathbf{a}\in \mathbb{R}^n \backslash\mathbf{0}$, $\alpha\in \mathbb{R}$, $\mathcal{H} = \{x\in \mathbb{R}^n : \mathbf{a}^T\mathbf{x} = \alpha\}$ and corresponding $\mathcal{H}^+$ s.t. $P \subseteq \mathcal{H}^+$ and $P \cap \mathcal{H} = \{\mathbf{\hat{x}}\}$. Do by contradiction: $\mathbf{\hat{x}}$ is feasible so for non-b.f. cannot be basic. So $\{\mathbf{a}^i : i\in I(\mathbf{\hat{x}})\}$ doesn't have $n$ lin-ind vecs. Then $\exists \mathbf{d}\in \mathbb{R}^n \backslash\mathbf{0}$ s.t. $(\mathbf{a}^{(i)})^T\mathbf{d} = \mathbf{0}$ for all $i\in I(\mathbf{\hat{x}})$. Take $\epsilon > 0$ and $\mathbf{\hat{x}} \pm\epsilon\mathbf{d}$. $(\mathbf{a}^{(i)})^T\mathbf{d} = \mathbf{0}$ for $i\in I(\mathbf{\hat{x}})$ implies $(\mathbf{a}^{(i)})^T(\mathbf{x} \pm\epsilon\mathbf{d}) = b_i$ for $i\in I(\mathbf{\hat{x}})$. For $i\in M_1\backslash I(\mathbf{\hat{x}})$ ($\in M_2\backslash I(\mathbf{\hat{x}})$) must have $(\mathbf{a}^{(i)})^T\mathbf{x} > (<) b_i$ so $(\mathbf{a}^{(i)})^T(\mathbf{\hat{x}} +\epsilon\mathbf{d}) \ge (\le)$ for $\epsilon$ small enough. So can find $\epsilon^* > 0$ s.t. $\mathbf{\hat{x}} \pm \epsilon^*\mathbf{d}$ both $\in \mathcal{P}$, apply $x\in \mathcal{P} \implies \mathbf{a}^T \mathbf{x} \ge \alpha$ for $\mathbf{a}^T\mathbf{d} >0$ and $<0$ for contradiction.
			\end{proof}
			
		\end{thm}
		
		
		
		{\tiny\textbf{\underline{Existence Of Vertices}}}
		
		
		\begin{dfn}[Lines]{dfn:line}{ }
			A polyhedron $\mathcal{P}$ contains a line if $\exists \mathbf{\hat{x}}\in \mathcal{P}$ and $\mathbf{d}\in \mathbb{R}^n \backslash \mathbf{0}$ s.t. $\mathbf{\hat{x}} + \lambda \mathbf{d}\in \mathcal{P}$ for all $\lambda\in \mathbb{R}$.
			\begin{itemize}
				\item $\mathcal{P}$ has a vertex iff it does not contain a line
				
%%%% UNFINISHED !!!!!!!!!!!!! PROBS NOT WORTH PUTTING IN THO
%				\vspace{-7pt}
%				\begin{proof}
%					$\implies$: By contrapositive, since $(\mathbf{a}^{(i)})^T \mathbf{\hat{x}}$ and $(\mathbf{a}^{(i)})^T (\mathbf{\hat{x}} + \lambda \mathbf{d})$ both satisfy constraints re $b_i$ we can get $(\mathbf{a}^{(i)})^T \mathbf{d} = \mathbf{0}$ for $i\in M_1 \cup M_2 \cup M_3$ so for all $\mathbf{x}\in \mathcal{P}$, $(\mathbf{a}^{(i)})^T\mathbf{d} = \mathbf{0}$ for $i\in I(\mathbf{x})$ so $\{a^{(i)} : i\in I(\mathbf{x})\}$ cannot have $n$ lin-ind vecs, so no $\mathbf{x}\in \mathcal{P}$ can be a vertex.\\
%					$\impliedby$: By last part, if $\nexists$ a line, $\nexists \mathbf{d}$ sending all $(\mathbf{a}^{(i)})^T$ to $\mathbf{0}$. Want $\mathbf{\hat{x}}\in \mathcal{P}$ s.t. $\{\mathbf{a}^{(i)} : i\in I(\mathbf{\hat{arg}})\}$ has n lin-ind vecs. If it doesn't for some $\mathbf{\hat{x}}\in \mathcal{P}$, a $\mathbf{d}$ sending all $(\mathbf{a}^{(i)})^T$ exists. Go on the line $\mathbf{\hat{x}} + \lambda\mathbf{d}$ to $\mathbf{x}^*$ where a new constraint $i^*$ binds, $(\mathbf{a}^{(i^*)})^T \mathbf{d} \ne 0$ and $I(\mathbf{x}^*) \supseteq I(\mathbf{\hat{x}})$, $(\mathbf{a}^{(i^*)})^T$ not a lin-comb of $						
%				\end{proof}
				\item Nonempty polytopes have vertices (since they cannot contain lines)
				\item Nonempty polyhedra have at most a finite number of vertices (if they have no lines, any $\mathbf{\hat{x}}\in \mathcal{P}$ has $I(\mathbf{\hat{x}}) \subseteq M_3 \cup J$ for $J \subseteq M_1 \cup M_2$, so finite number of sets of binding constraints. Only a subset of those have $n$ lin-ind $\mathbf{a}^{(i)}$, if $\mathbf{x}_1$ and $\mathbf{x}_2$ have same $I$ and $I$ gives $n$ lin-ind $\mathbf{a}^{(i)}$, then $(\mathbf{a}^{(i)})^T(\mathbf{x}_1 - \mathbf{x})_2 = \mathbf{0}$ for $i\in I$.)
				\item General convex sets can have an infinite number of vertices (e.g. a circle)
			\end{itemize}
			
		\end{dfn}
		
		
		{\tiny\textbf{\underline{Graphical Solution Method}}}
		
		
		\begin{rem}[Graphical Solution Method]{rem:graphicalsolution}{ }
			Only works in 1D and 2D really, if Miguel asks us to draw in 3D there will be tears. Just draw the feasible region and $-\mathbf{c}$ then go as far in the $-\mathbf{c}$ direction as possible. Remember that an optimal solution will always occur at some vertex if one exists! 
			
		\end{rem}
		
		
		
		
		UP TO PAGE 44 - START OF LECTURE 10
		
		
\end{multicols*}	
	
\end{document}
